% ------------------------------------------------------------%
% 2021-06-03
% 2015-2026 - Emerson Ribeiro de Mello - mello@ifsc.edu.br
% ------------------------------------------------------------%
\documentclass[11pt,addpoints]{ifscexam}
\usepackage{ifscutils}

% ---------------------------------------------- %
% Para apresentar um formato de questão diferente
% ---------------------------------------------- %
% \formatoquestao
% ---------------------------------------------- %

% ---------------------------------------------- %
% Para imprimir o gabarito
% ---------------------------------------------- %
% \printanswers

% Em uma caixa com o fundo preenchido
\shadedsolutions
% Ou somente uma caixa sem fundo
%\framedsolutions
% ---------------------------------------------- %


\begin{document}


\cabecalhoCurso%
{Engenharia de Telecomunicações}% Nome do curso
{Nome da disciplina}% Nome da disciplina
{Prof. Emerson Ribeiro de Mello}% Nome do professor
{\url{mello@ifsc.edu.br}}% E-mail do professor
{\textbf{Avaliação 1}}% Título da avaliação
{24/11/2020}% Data da avaliação


% ---------------------------------------------- %
% Para imprimir espaço para nome do aluno
% ---------------------------------------------- %
\nomealuno{\rule{.58\linewidth}{.4pt}}

% Ou comente a linha acima e preencha com a matrícula e nome
% como apresentado abaixo

% \nomealuno[20201130123]{Nome Completo do Aluno}
% ---------------------------------------------- %

\begin{questions}

\question[1]
Qual a cor do céu?

\begin{solutionorlines}[2cm]
Geralmente é azul.
\end{solutionorlines}


\question[1]
Calcule $\displaystyle\int_0^1 x \, dx$.

\begin{solutionorbox}[2.5cm]
$\displaystyle\tfrac{x^2}{2}\Big|^1_0 = \tfrac{1^2}{2} - \tfrac{0^2}{2} = \tfrac{1}{2}$
\end{solutionorbox}

\question[1]
Quantos meses possuem no máximo 30 dias?
\begin{solutionorlines}[2cm]
Somente 4. São eles: abril, junho, setembro e novembro.
\end{solutionorlines}


\question[1]
Sobre o sistema operacional Linux
\begin{parts}
    \part
    Liste três aplicativos utilitários que vem por padrão
    \begin{solutionorlines}[1cm]
    ls, cd e mv.
    \end{solutionorlines}
    \part
    Liste o nome de três distribuições
    \begin{solutionorlines}[1cm]
    Debian, Redhat e Ubuntu. 
    \end{solutionorlines}
    \part 
    Explique a diferença entre um software aplicativo e um aplicativo utilitário
    \begin{solutionorlines}[3cm]
    Um software aplicativo é um programa que tem uma função específica, como um editor de texto ou um navegador de internet. Já um aplicativo utilitário é um programa que tem a função de auxiliar na administração do sistema operacional, como um gerenciador de arquivos ou um monitor de recursos.
    \end{solutionorlines}
\end{parts}

\question[2]
O sistema operacional Linux foi criado por \fillin[Linus Torvalds][3.5cm].

\question[1]
Quais dos meses abaixo pode ter um número de dias diferente em anos bissextos?
\begin{choices}
    \choice Janeiro
    \correctchoice Fevereiro
    \choice Abril
    \choice Junho
    \choice Setembro
\end{choices}

\question[1]
Marque todas as opções abaixo que são nomes de sistemas operacionais:
\begin{checkboxes}
    \correctchoice Linux
    \correctchoice Windows
    \choice LibreOffice
    \choice Bloco de notas
\end{checkboxes}


\fullwidth{\Large \textbf{Verdadeiro ou Falso}}

\question[1] \tf[V] O FreeBSD é um sistema operacional.
\question[1] \tf[F] O OpenBSD é baseado no sistema operacional Linux.

\end{questions}



% ---------------------------------------------- %
% Para imprimir tabela com as notas
% ---------------------------------------------- %
\vfill
\rule{10cm}{1px}

\vspace{0.5cm}
\gradetable[h]
% ---------------------------------------------- %
\end{document}